\selectlanguage{english}
\chapter*{Abstract}
\addcontentsline{toc}{chapter}{Abstract}

This thesis investigates the impact of confidence function design on the quality of recommendations generated by the implicit ALS model. In the method proposed by Hu, Koren, and Volinsky, observations are incorporated into the model with weights determined as a function of the intensity of user--item interactions. In research practice, however, the choice of this function is often treated as an implicit assumption rather than being systematically validated. The objective of this thesis was to examine the extent to which different transformations of interaction intensity affect recommendation quality and to compare this impact with that of the choice of the recommendation model itself.

The study was conducted using real-world data from the Steam platform, comprising users' game libraries together with the recorded cumulative playtime for each title. Five confidence function variants were compared: linear, logarithmic, logarithmic with saturation, logarithmic with offset, and a constant function that completely ignores playtime. These variants were evaluated alongside the Bayesian Personalised Ranking model and a popularity-based recommender. The hyperparameters of each model were tuned independently using Bayesian optimization. Recommendation quality was evaluated using ranking accuracy metrics as well as beyond-accuracy measures, including catalog coverage, intra-list diversity, and popularity bias. The study also included an analysis of result stability with respect to random data splits and model initialization, statistical significance testing with effect size estimation, and a sensitivity analysis of the confidence function parameters.

The results indicate that weighting confidence according to playtime does not improve recommendation quality. The variant with the constant confidence function achieved the highest accuracy, outperforming the logarithmic variant by 6.2\% and the linear variant by 8.1\%. At the same time, the performance gap between the implicit ALS models and the BPR model was several times larger than the entire range of differences observed within the implicit ALS family, while the BPR model achieved more than twice the catalog coverage. Consequently, the choice of the recommendation model primarily determines the trade-off between the different dimensions of recommendation quality.

The findings suggest that, for the dataset considered, the mere presence of a game in a user's library constitutes a sufficient preference signal, whereas attempts to extract additional information from interaction duration provide no added value beyond increasing the overall confidence weights assigned to observed interactions. These conclusions apply specifically to the interaction signal defined as the cumulative playtime accumulated over the entire ownership period of a game.

\vspace{1em}
\noindent\textbf{Keywords:} recommender systems, implicit feedback, matrix factorization, implicit ALS, confidence function, Steam platform
\selectlanguage{polish}
